\section{Data Curation and Ingestion Pipelines}
\label{sec:pipelines}

To support deterministic theorem proving, Omni-IPS requires structured knowledge representation across all supported domains. We designed separate, robust ingestion pipelines in \texttt{data\_pipelines/} to crawl, curate, and format domain-specific rules and facts prior to seeding our dual-database storage. The workflow uses strict validation schemas and separates structural definitions into clear domains.

\subsection{Ingestion Pipelines Workflow}
The ingestion process is unified across the Chemistry, Geometry, and Algebra domains, executing the following pipeline steps:
\begin{enumerate}
    \item \textbf{Schema Definition}: Establishing rigid JSON formats for facts and rules.
    \item \textbf{Domain Curation}: Aggregating historical axioms, chemical equations, and algebraic properties.
    \item \textbf{Database Ingest}: Detaching legacy nodes and writing atomic properties and logical connections to Neo4j.
    \item \textbf{Semantic Embedding}: Encoding mathematical and textual descriptions into high-dimensional vector representations stored in Qdrant.
\end{enumerate}

\subsection{Chemistry Ingestion Pipeline}
The chemistry pipeline in \texttt{ingest\_chemistry.py} processes chemical reactions, elements, and stoichiometric compounds. Facts are modeled as chemical species (atoms, molecules, or ions) and rules represent chemical reactions. 
For example, the single-displacement reaction of sodium in water is modeled as a rule:
\begin{itemize}
    \item \textbf{Antecedents}: $\text{Na}$, $\text{H}_2\text{O}$
    \item \textbf{Consequents}: $\text{NaOH}$, $\text{H}_2$
    \item \textbf{Stoichiometry Description}: $2\text{Na} + 2\text{H}_2\text{O} \to 2\text{NaOH} + \text{H}_2$
\end{itemize}

Chemical properties (e.g., aggregation states, molecular weights, safety classifications) are stored in the node attributes. We curate common reaction archetypes including:
\begin{table}[h]
    \centering
    \caption{Ingested Chemical Reaction Archetypes in Omni-IPS}
    \label{tab:chem_reactions}
    \begin{tabular}{llll}
        \toprule
        \textbf{Reaction Type} & \textbf{Inputs} & \textbf{Outputs} & \textbf{Example Equation} \\
        \midrule
        Synthesis & $\text{H}_2, \text{O}_2$ & $\text{H}_2\text{O}$ & $2\text{H}_2 + \text{O}_2 \to 2\text{H}_2\text{O}$ \\
        Decomposition & $\text{CaCO}_3$ & $\text{CaO}, \text{CO}_2$ & $\text{CaCO}_3 \to \text{CaO} + \text{CO}_2$ \\
        Neutralization & $\text{NaOH}, \text{HCl}$ & $\text{NaCl}, \text{H}_2\text{O}$ & $\text{NaOH} + \text{HCl} \to \text{NaCl} + \text{H}_2\text{O}$ \\
        Precipitation & $\text{BaCl}_2, \text{Na}_2\text{SO}_4$ & $\text{BaSO}_4, \text{NaCl}$ & $\text{BaCl}_2 + \text{Na}_2\text{SO}_4 \to \text{BaSO}_4 \downarrow + 2\text{NaCl}$ \\
        \bottomrule
    \end{tabular}
\end{table}

\subsection{Geometry Ingestion Pipeline}
The geometry pipeline in \texttt{ingest\_geometry.py} models Euclidean flat geometry axioms, triangles, angles, lines, and congruence. The main challenge in Euclidean geometry is representing spatial relationships symbolically. Facts are represented as geometric predicates:
\begin{equation}
    R(A_1, A_2, \dots, A_n)
\end{equation}
where $R$ is a geometric relation (e.g., \texttt{RightTriangle}, \texttt{Parallel}, \texttt{Perpendicular}) and $A_i$ represent spatial vertices, lines, or angles. Rules represent Euclidean theorems. For instance, the transitivity of parallel lines:
\begin{equation}
    \texttt{Parallel}(a,b) \land \texttt{Parallel}(b,c) \implies \texttt{Parallel}(a,c)
\end{equation}

\subsection{Algebra Ingestion Pipeline}
The algebra pipeline in \texttt{ingest\_algebra.py} handles basic symbolic equations and linear system manipulation rules. Unlike geometry, algebra relies heavily on equation rewrites and variable substitution. Facts represent equation states, such as \texttt{Equation(x+2, 5)} or simple constant mappings. Rules represent algebraic properties of equality (e.g., subtracting or adding constants to both sides of an equation) and quadratic properties:
\begin{equation}
    a \cdot x^2 + b \cdot x + c = 0 \implies x = \frac{-b \pm \sqrt{b^2 - 4ac}}{2a}
\end{equation}
This allows the solver to logically transition from initial inputs to solved states by matching exact algebraic rewrite steps.
